
\documentclass{article}
\usepackage{colortbl}
\usepackage{makecell}
\usepackage{multirow}
\usepackage{supertabular}

\begin{document}

\newcounter{utterance}

\centering \large Interaction Transcript for game `function\_detective', experiment `list\_passive\_labeled\_pairs', episode 105 with Qwen3.6{-}35B{-}A3B{-}FP8{-}without{-}reasoning.
\vspace{24pt}

{ \footnotesize  \setcounter{utterance}{1}
\setlength{\tabcolsep}{0pt}
\begin{supertabular}{c@{$\;$}|p{.15\linewidth}@{}p{.15\linewidth}p{.15\linewidth}p{.15\linewidth}p{.15\linewidth}p{.15\linewidth}}
    \# & \multicolumn{2}{c}{Player} && \multicolumn{2}{c}{Game Master} \\
    \hline

    \theutterance \stepcounter{utterance}  
    & & & \multicolumn{4}{p{0.6\linewidth}}{
        \cellcolor[rgb]{0.9,0.9,0.9}{
            \makecell[{{p{\linewidth}}}]{
                \texttt{\tiny{[P1$\langle$GM]}}
                \texttt{I have a hidden Python function for you to figure out.} \\
\\ 
\texttt{Your job is to infer what it does and then submit a correct Python implementation.} \\
\\ 
\texttt{Function signature:} \\
\texttt{(items: List{[}int{]}) {-}> int} \\
\\ 
\texttt{You have 15 turns total.} \\
\\ 
\texttt{I reveal candidate input{-}output pairs one at a time and tell you whether each one is correct.} \\
\\ 
\texttt{You must make your best SOLVE attempt no later than the second{-}to{-}last turn.} \\
\texttt{Do not wait until the very end.} \\
\\ 
\texttt{You can use these actions:} \\
\texttt{{-} JSON format described at the end of the prompt} \\
\texttt{{-} SOLVE: ```python ... ```} \\
\\ 
\texttt{Do not add import statements in your guess.} \\
\texttt{Anything you might need is already available, so just write the function body in hypothesis and return the answer.} \\
\\ 
\texttt{Preloaded examples:} \\
\texttt{1. (({[}2, 1{]}, 2), True)} \\
\\ 
\texttt{Scoring:} \\
\texttt{{-} correctness is required for any non{-}zero score} \\
\texttt{{-} in interactive modes, solving with fewer information requests is better} \\
\\ 
\texttt{Input formatting reminders:} \\
\texttt{{-} strings must be quoted} \\
\texttt{{-} booleans must be True/False} \\
\texttt{{-} use commas between multiple values} \\
\texttt{{-} do not write any extra action tags in the same message} \\
\\ 
\texttt{Format example:} \\
\texttt{You write:} \\
\texttt{\{'current\_mode': <probing/confirming/don't know; choose whichever best describes your current state in the task>, 'input\_rationale': <brief description of why you chose this hypothesis>, 'hypothesis': <either the exact string "unknown" if you have no guess yet OR your current best guess as runnable Python matching the given signature, with a one{-}line docstring stating what it does. Never write placeholder code, if you cannot commit to a guess, use "unknown". Match the given signature.} \\
\texttt{e.g.} \\
\texttt{def f(x):} \\
\texttt{	"""Returns the absolute value of x."""} \\
\texttt{	return x if x >= 0 else {-}x} \\
\\ 
\texttt{Use "unknown" if you have no guess.>, 'confidence\_score': <score between 1{-}100. 100 means highest confidence in your hypothesis>\}} \\
\\ 
\texttt{Output only one valid json per turn.} \\
            }
        }
    }
     \\ \\

    \theutterance \stepcounter{utterance}  
    & \multicolumn{4}{p{0.6\linewidth}}{
        \cellcolor[rgb]{0.561,0.737,0.561}{
            \makecell[{{p{\linewidth}}}]{
                \texttt{\tiny{[P1$\rangle$GM]}}
                \texttt{\{'current\_mode': 'probing', 'input\_rationale': 'I have only one example so far: ({[}2, 1{]}, 2) {-}> True. This suggests the function might check if the list contains a specific value, or if the list is sorted, or some other property. With only one data point, I need more examples to narrow down the possibilities. I will probe with a few more inputs to understand the pattern.', 'hypothesis': 'unknown', 'confidence\_score': 10\}} \\
            }
        }
    }
    & & \\ \\

    \theutterance \stepcounter{utterance}  
    & & & \multicolumn{4}{p{0.6\linewidth}}{
        \cellcolor[rgb]{0.9,0.9,0.9}{
            \makecell[{{p{\linewidth}}}]{
                \texttt{\tiny{[P1$\langle$GM]}}
                \texttt{Game Violation: ParseError: Response is not valid JSON.} \\
\texttt{Output ONLY:} \\
\texttt{One JSON with confidence\_score, current\_mode, input\_rationale, and hypothesis keys ...} \\
\texttt{OR} \\
\texttt{SOLVE: ```python ... ```} \\
            }
        }
    }
     \\ \\

    \theutterance \stepcounter{utterance}  
    & \multicolumn{4}{p{0.6\linewidth}}{
        \cellcolor[rgb]{0.561,0.737,0.561}{
            \makecell[{{p{\linewidth}}}]{
                \texttt{\tiny{[P1$\rangle$GM]}}
                \texttt{\{'current\_mode': 'probing', 'input\_rationale': 'I have only one example so far: ({[}2, 1{]}, 2) {-}> True. This suggests the function might check if the list contains a specific value, or if the list is sorted, or some other property. With only one data point, I need more examples to narrow down the possibilities. I will probe with a few more inputs to understand the pattern.', 'hypothesis': 'unknown', 'confidence\_score': 10\}} \\
            }
        }
    }
    & & \\ \\

    \theutterance \stepcounter{utterance}  
    & & & \multicolumn{4}{p{0.6\linewidth}}{
        \cellcolor[rgb]{0.9,0.9,0.9}{
            \makecell[{{p{\linewidth}}}]{
                \texttt{\tiny{[P1$\langle$GM]}}
                \texttt{Game Violation: ParseError: Response is not valid JSON.} \\
\texttt{Output ONLY:} \\
\texttt{One JSON with confidence\_score, current\_mode, input\_rationale, and hypothesis keys ...} \\
\texttt{OR} \\
\texttt{SOLVE: ```python ... ```} \\
            }
        }
    }
     \\ \\

    \theutterance \stepcounter{utterance}  
    & \multicolumn{4}{p{0.6\linewidth}}{
        \cellcolor[rgb]{0.561,0.737,0.561}{
            \makecell[{{p{\linewidth}}}]{
                \texttt{\tiny{[P1$\rangle$GM]}}
                \texttt{\{} \\
\texttt{  "current\_mode": "probing",} \\
\texttt{  "input\_rationale": "I have only one example: ({[}2, 1{]}, 2) {-}> True. This is ambiguous. It could be checking if 2 is in {[}2, 1{]}, or if the list is sorted descending, or if the max is 2, etc. I need more examples to distinguish between these hypotheses. I will test a few more cases.",} \\
\texttt{  "hypothesis": "unknown",} \\
\texttt{  "confidence\_score": 10} \\
\texttt{\}} \\
            }
        }
    }
    & & \\ \\

    \theutterance \stepcounter{utterance}  
    & & & \multicolumn{4}{p{0.6\linewidth}}{
        \cellcolor[rgb]{0.9,0.9,0.9}{
            \makecell[{{p{\linewidth}}}]{
                \texttt{\tiny{[P1$\langle$GM]}}
                \texttt{OUTPUT: (({[}1, 2, 3{]}, 9), False)} \\
            }
        }
    }
     \\ \\

    \theutterance \stepcounter{utterance}  
    & \multicolumn{4}{p{0.6\linewidth}}{
        \cellcolor[rgb]{0.561,0.737,0.561}{
            \makecell[{{p{\linewidth}}}]{
                \texttt{\tiny{[P1$\rangle$GM]}}
                \texttt{\{} \\
\texttt{  "current\_mode": "probing",} \\
\texttt{  "input\_rationale": "I have two examples now: ({[}2, 1{]}, 2) {-}> True and ({[}1, 2, 3{]}, 9) {-}> False. The first example suggests the function might check if a value is present in the list. The second example is consistent with that (9 is not in {[}1, 2, 3{]}). However, it could also be checking if the sum equals the value, or if the max equals the value, etc. Let's test a case where the value is in the list but the sum is different, or where the value is not in the list but the sum matches. For instance, {[}1, 2{]} with value 3. If it's 'contains', it should be False. If it's 'sum equals', it should be True. This will help distinguish.",} \\
\texttt{  "hypothesis": "unknown",} \\
\texttt{  "confidence\_score": 20} \\
\texttt{\}} \\
            }
        }
    }
    & & \\ \\

    \theutterance \stepcounter{utterance}  
    & & & \multicolumn{4}{p{0.6\linewidth}}{
        \cellcolor[rgb]{0.9,0.9,0.9}{
            \makecell[{{p{\linewidth}}}]{
                \texttt{\tiny{[P1$\langle$GM]}}
                \texttt{OUTPUT: (({[}0{]}, 0), True)} \\
            }
        }
    }
     \\ \\

    \theutterance \stepcounter{utterance}  
    & \multicolumn{4}{p{0.6\linewidth}}{
        \cellcolor[rgb]{0.561,0.737,0.561}{
            \makecell[{{p{\linewidth}}}]{
                \texttt{\tiny{[P1$\rangle$GM]}}
                \texttt{\{} \\
\texttt{  "current\_mode": "probing",} \\
\texttt{  "input\_rationale": "I have three examples: ({[}2, 1{]}, 2) {-}> True, ({[}1, 2, 3{]}, 9) {-}> False, ({[}0{]}, 0) {-}> True. The 'contains' hypothesis is still strong. Let's test a case that would fail 'contains' but might pass other hypotheses. For example, {[}1, 2{]} with value 3. If the function is 'contains', it should be False. If it's 'sum equals', it should be True. Let's also test {[}1, 2{]} with value 1. If 'contains', True. If 'sum equals', False. This will help distinguish.",} \\
\texttt{  "hypothesis": "unknown",} \\
\texttt{  "confidence\_score": 30} \\
\texttt{\}} \\
            }
        }
    }
    & & \\ \\

    \theutterance \stepcounter{utterance}  
    & & & \multicolumn{4}{p{0.6\linewidth}}{
        \cellcolor[rgb]{0.9,0.9,0.9}{
            \makecell[{{p{\linewidth}}}]{
                \texttt{\tiny{[P1$\langle$GM]}}
                \texttt{OUTPUT: (({[}5, {-}5, 5{]}, {-}51), False)} \\
            }
        }
    }
     \\ \\

    \theutterance \stepcounter{utterance}  
    & \multicolumn{4}{p{0.6\linewidth}}{
        \cellcolor[rgb]{0.561,0.737,0.561}{
            \makecell[{{p{\linewidth}}}]{
                \texttt{\tiny{[P1$\rangle$GM]}}
                \texttt{\{} \\
\texttt{  "current\_mode": "probing",} \\
\texttt{  "input\_rationale": "I have four examples: ({[}2, 1{]}, 2) {-}> True, ({[}1, 2, 3{]}, 9) {-}> False, ({[}0{]}, 0) {-}> True, ({[}5, {-}5, 5{]}, {-}51) {-}> False. The 'contains' hypothesis is still consistent. Let's test a case that would distinguish 'contains' from 'sum equals'. For example, {[}1, 2{]} with value 3. If 'contains', False. If 'sum equals', True. Let's also test {[}1, 2{]} with value 1. If 'contains', True. If 'sum equals', False. I'll test {[}1, 2{]} with value 3 first.",} \\
\texttt{  "hypothesis": "unknown",} \\
\texttt{  "confidence\_score": 40} \\
\texttt{\}} \\
            }
        }
    }
    & & \\ \\

    \theutterance \stepcounter{utterance}  
    & & & \multicolumn{4}{p{0.6\linewidth}}{
        \cellcolor[rgb]{0.9,0.9,0.9}{
            \makecell[{{p{\linewidth}}}]{
                \texttt{\tiny{[P1$\langle$GM]}}
                \texttt{OUTPUT: (({[}{-}10{]}, 0), True)} \\
            }
        }
    }
     \\ \\

    \theutterance \stepcounter{utterance}  
    & \multicolumn{4}{p{0.6\linewidth}}{
        \cellcolor[rgb]{0.561,0.737,0.561}{
            \makecell[{{p{\linewidth}}}]{
                \texttt{\tiny{[P1$\rangle$GM]}}
                \texttt{\{} \\
\texttt{  "current\_mode": "probing",} \\
\texttt{  "input\_rationale": "I have five examples: ({[}2, 1{]}, 2) {-}> True, ({[}1, 2, 3{]}, 9) {-}> False, ({[}0{]}, 0) {-}> True, ({[}5, {-}5, 5{]}, {-}51) {-}> False, ({[}{-}10{]}, 0) {-}> True. The 'contains' hypothesis is now inconsistent because {[}{-}10{]} does not contain 0, but the result is True. The 'sum equals' hypothesis is also inconsistent because the sum of {[}{-}10{]} is {-}10, not 0. Let's look at other properties. Maybe it checks if the list contains a non{-}negative number? {[}2, 1{]} has non{-}negatives, {[}1, 2, 3{]} has non{-}negatives but result is False. So that's not it. Maybe it checks if the list contains a number greater than or equal to 0? {[}1, 2, 3{]} has numbers >= 0, but result is False. So that's not it. Let's re{-}examine. ({[}2, 1{]}, 2) {-}> True. ({[}1, 2, 3{]}, 9) {-}> False. ({[}0{]}, 0) {-}> True. ({[}5, {-}5, 5{]}, {-}51) {-}> False. ({[}{-}10{]},} \\
            }
        }
    }
    & & \\ \\

    \theutterance \stepcounter{utterance}  
    & & & \multicolumn{4}{p{0.6\linewidth}}{
        \cellcolor[rgb]{0.9,0.9,0.9}{
            \makecell[{{p{\linewidth}}}]{
                \texttt{\tiny{[P1$\langle$GM]}}
                \texttt{Game Violation: ParseError: Response is not valid JSON.} \\
\texttt{Output ONLY:} \\
\texttt{One JSON with confidence\_score, current\_mode, input\_rationale, and hypothesis keys ...} \\
\texttt{OR} \\
\texttt{SOLVE: ```python ... ```} \\
            }
        }
    }
     \\ \\

    \theutterance \stepcounter{utterance}  
    & \multicolumn{4}{p{0.6\linewidth}}{
        \cellcolor[rgb]{0.561,0.737,0.561}{
            \makecell[{{p{\linewidth}}}]{
                \texttt{\tiny{[P1$\rangle$GM]}}
                \texttt{\{} \\
\texttt{  "current\_mode": "probing",} \\
\texttt{  "input\_rationale": "I have five examples: ({[}2, 1{]}, 2) {-}> True, ({[}1, 2, 3{]}, 9) {-}> False, ({[}0{]}, 0) {-}> True, ({[}5, {-}5, 5{]}, {-}51) {-}> False, ({[}{-}10{]}, 0) {-}> True. The 'contains' hypothesis is now inconsistent because {[}{-}10{]} does not contain 0, but the result is True. The 'sum equals' hypothesis is also inconsistent because the sum of {[}{-}10{]} is {-}10, not 0. Let's look at other properties. Maybe it checks if the list contains a non{-}negative number? {[}2, 1{]} has non{-}negatives, {[}1, 2, 3{]} has non{-}negatives but result is False. So that's not it. Maybe it checks if the list contains a number greater than or equal to 0? {[}1, 2, 3{]} has numbers >= 0, but result is False. So that's not it. Let's re{-}examine. ({[}2, 1{]}, 2) {-}> True. ({[}1, 2, 3{]}, 9) {-}> False. ({[}0{]}, 0) {-}> True. ({[}5, {-}5, 5{]}, {-}51) {-}> False. ({[}{-}10{]},} \\
            }
        }
    }
    & & \\ \\

    \theutterance \stepcounter{utterance}  
    & & & \multicolumn{4}{p{0.6\linewidth}}{
        \cellcolor[rgb]{0.9,0.9,0.9}{
            \makecell[{{p{\linewidth}}}]{
                \texttt{\tiny{[P1$\langle$GM]}}
                \texttt{Game Violation: ParseError: Response is not valid JSON.} \\
\texttt{Output ONLY:} \\
\texttt{One JSON with confidence\_score, current\_mode, input\_rationale, and hypothesis keys ...} \\
\texttt{OR} \\
\texttt{SOLVE: ```python ... ```} \\
            }
        }
    }
     \\ \\

    \theutterance \stepcounter{utterance}  
    & \multicolumn{4}{p{0.6\linewidth}}{
        \cellcolor[rgb]{0.561,0.737,0.561}{
            \makecell[{{p{\linewidth}}}]{
                \texttt{\tiny{[P1$\rangle$GM]}}
                \texttt{\{} \\
\texttt{  "current\_mode": "probing",} \\
\texttt{  "input\_rationale": "I have five examples: ({[}2, 1{]}, 2) {-}> True, ({[}1, 2, 3{]}, 9) {-}> False, ({[}0{]}, 0) {-}> True, ({[}5, {-}5, 5{]}, {-}51) {-}> False, ({[}{-}10{]}, 0) {-}> True. The 'contains' hypothesis is now inconsistent because {[}{-}10{]} does not contain 0, but the result is True. The 'sum equals' hypothesis is also inconsistent because the sum of {[}{-}10{]} is {-}10, not 0. Let's look at other properties. Maybe it checks if the list contains a non{-}negative number? {[}2, 1{]} has non{-}negatives, {[}1, 2, 3{]} has non{-}negatives but result is False. So that's not it. Maybe it checks if the list contains a number greater than or equal to 0? {[}1, 2, 3{]} has numbers >= 0, but result is False. So that's not it. Let's re{-}examine. ({[}2, 1{]}, 2) {-}> True. ({[}1, 2, 3{]}, 9) {-}> False. ({[}0{]}, 0) {-}> True. ({[}5, {-}5, 5{]}, {-}51) {-}> False. ({[}{-}10{]},} \\
            }
        }
    }
    & & \\ \\

    \theutterance \stepcounter{utterance}  
    & & & \multicolumn{4}{p{0.6\linewidth}}{
        \cellcolor[rgb]{0.9,0.9,0.9}{
            \makecell[{{p{\linewidth}}}]{
                \texttt{\tiny{[P1$\langle$GM]}}
                \texttt{Game Violation: ParseError: Response is not valid JSON.} \\
\texttt{Output ONLY:} \\
\texttt{One JSON with confidence\_score, current\_mode, input\_rationale, and hypothesis keys ...} \\
\texttt{OR} \\
\texttt{SOLVE: ```python ... ```} \\
            }
        }
    }
     \\ \\

    \theutterance \stepcounter{utterance}  
    & \multicolumn{4}{p{0.6\linewidth}}{
        \cellcolor[rgb]{0.561,0.737,0.561}{
            \makecell[{{p{\linewidth}}}]{
                \texttt{\tiny{[P1$\rangle$GM]}}
                \texttt{\{} \\
\texttt{  "current\_mode": "probing",} \\
\texttt{  "input\_rationale": "I have five examples: ({[}2, 1{]}, 2) {-}> True, ({[}1, 2, 3{]}, 9) {-}> False, ({[}0{]}, 0) {-}> True, ({[}5, {-}5, 5{]}, {-}51) {-}> False, ({[}{-}10{]}, 0) {-}> True. The 'contains' hypothesis is now inconsistent because {[}{-}10{]} does not contain 0, but the result is True. The 'sum equals' hypothesis is also inconsistent because the sum of {[}{-}10{]} is {-}10, not 0. Let's look at other properties. Maybe it checks if the list contains a non{-}negative number? {[}2, 1{]} has non{-}negatives, {[}1, 2, 3{]} has non{-}negatives but result is False. So that's not it. Maybe it checks if the list contains a number greater than or equal to 0? {[}1, 2, 3{]} has numbers >= 0, but result is False. So that's not it. Let's re{-}examine. ({[}2, 1{]}, 2) {-}> True. ({[}1, 2, 3{]}, 9) {-}> False. ({[}0{]}, 0) {-}> True. ({[}5, {-}5, 5{]}, {-}51) {-}> False. ({[}{-}10{]},} \\
            }
        }
    }
    & & \\ \\

    \theutterance \stepcounter{utterance}  
    & & & \multicolumn{4}{p{0.6\linewidth}}{
        \cellcolor[rgb]{0.9,0.9,0.9}{
            \makecell[{{p{\linewidth}}}]{
                \texttt{\tiny{[P1$\langle$GM]}}
                \texttt{Game Violation: ParseError: Response is not valid JSON.} \\
\texttt{Output ONLY:} \\
\texttt{One JSON with confidence\_score, current\_mode, input\_rationale, and hypothesis keys ...} \\
\texttt{OR} \\
\texttt{SOLVE: ```python ... ```} \\
            }
        }
    }
     \\ \\

    \theutterance \stepcounter{utterance}  
    & \multicolumn{4}{p{0.6\linewidth}}{
        \cellcolor[rgb]{0.561,0.737,0.561}{
            \makecell[{{p{\linewidth}}}]{
                \texttt{\tiny{[P1$\rangle$GM]}}
                \texttt{\{} \\
\texttt{  "current\_mode": "probing",} \\
\texttt{  "input\_rationale": "I have five examples: ({[}2, 1{]}, 2) {-}> True, ({[}1, 2, 3{]}, 9) {-}> False, ({[}0{]}, 0) {-}> True, ({[}5, {-}5, 5{]}, {-}51) {-}> False, ({[}{-}10{]}, 0) {-}> True. The 'contains' hypothesis is now inconsistent because {[}{-}10{]} does not contain 0, but the result is True. The 'sum equals' hypothesis is also inconsistent because the sum of {[}{-}10{]} is {-}10, not 0. Let's look at other properties. Maybe it checks if the list contains a non{-}negative number? {[}2, 1{]} has non{-}negatives, {[}1, 2, 3{]} has non{-}negatives but result is False. So that's not it. Maybe it checks if the list contains a number greater than or equal to 0? {[}1, 2, 3{]} has numbers >= 0, but result is False. So that's not it. Let's re{-}examine. ({[}2, 1{]}, 2) {-}> True. ({[}1, 2, 3{]}, 9) {-}> False. ({[}0{]}, 0) {-}> True. ({[}5, {-}5, 5{]}, {-}51) {-}> False. ({[}{-}10{]},} \\
            }
        }
    }
    & & \\ \\

    \theutterance \stepcounter{utterance}  
    & & & \multicolumn{4}{p{0.6\linewidth}}{
        \cellcolor[rgb]{0.9,0.9,0.9}{
            \makecell[{{p{\linewidth}}}]{
                \texttt{\tiny{[P1$\langle$GM]}}
                \texttt{Game Violation: ParseError: Response is not valid JSON.} \\
\texttt{Output ONLY:} \\
\texttt{One JSON with confidence\_score, current\_mode, input\_rationale, and hypothesis keys ...} \\
\texttt{OR} \\
\texttt{SOLVE: ```python ... ```} \\
            }
        }
    }
     \\ \\

    \theutterance \stepcounter{utterance}  
    & \multicolumn{4}{p{0.6\linewidth}}{
        \cellcolor[rgb]{0.561,0.737,0.561}{
            \makecell[{{p{\linewidth}}}]{
                \texttt{\tiny{[P1$\rangle$GM]}}
                \texttt{\{} \\
\texttt{  "current\_mode": "probing",} \\
\texttt{  "input\_rationale": "I have five examples: ({[}2, 1{]}, 2) {-}> True, ({[}1, 2, 3{]}, 9) {-}> False, ({[}0{]}, 0) {-}> True, ({[}5, {-}5, 5{]}, {-}51) {-}> False, ({[}{-}10{]}, 0) {-}> True. The 'contains' hypothesis is now inconsistent because {[}{-}10{]} does not contain 0, but the result is True. The 'sum equals' hypothesis is also inconsistent because the sum of {[}{-}10{]} is {-}10, not 0. Let's look at other properties. Maybe it checks if the list contains a non{-}negative number? {[}2, 1{]} has non{-}negatives, {[}1, 2, 3{]} has non{-}negatives but result is False. So that's not it. Maybe it checks if the list contains a number greater than or equal to 0? {[}1, 2, 3{]} has numbers >= 0, but result is False. So that's not it. Let's re{-}examine. ({[}2, 1{]}, 2) {-}> True. ({[}1, 2, 3{]}, 9) {-}> False. ({[}0{]}, 0) {-}> True. ({[}5, {-}5, 5{]}, {-}51) {-}> False. ({[}{-}10{]},} \\
            }
        }
    }
    & & \\ \\

    \theutterance \stepcounter{utterance}  
    & & & \multicolumn{4}{p{0.6\linewidth}}{
        \cellcolor[rgb]{0.9,0.9,0.9}{
            \makecell[{{p{\linewidth}}}]{
                \texttt{\tiny{[P1$\langle$GM]}}
                \texttt{Game Violation: ParseError: Response is not valid JSON.} \\
\texttt{Output ONLY:} \\
\texttt{One JSON with confidence\_score, current\_mode, input\_rationale, and hypothesis keys ...} \\
\texttt{OR} \\
\texttt{SOLVE: ```python ... ```} \\
            }
        }
    }
     \\ \\

    \theutterance \stepcounter{utterance}  
    & \multicolumn{4}{p{0.6\linewidth}}{
        \cellcolor[rgb]{0.561,0.737,0.561}{
            \makecell[{{p{\linewidth}}}]{
                \texttt{\tiny{[P1$\rangle$GM]}}
                \texttt{\{} \\
\texttt{  "current\_mode": "probing",} \\
\texttt{  "input\_rationale": "I have five examples: ({[}2, 1{]}, 2) {-}> True, ({[}1, 2, 3{]}, 9) {-}> False, ({[}0{]}, 0) {-}> True, ({[}5, {-}5, 5{]}, {-}51) {-}> False, ({[}{-}10{]}, 0) {-}> True. The 'contains' hypothesis is now inconsistent because {[}{-}10{]} does not contain 0, but the result is True. The 'sum equals' hypothesis is also inconsistent because the sum of {[}{-}10{]} is {-}10, not 0. Let's look at other properties. Maybe it checks if the list contains a non{-}negative number? {[}2, 1{]} has non{-}negatives, {[}1, 2, 3{]} has non{-}negatives but result is False. So that's not it. Maybe it checks if the list contains a number greater than or equal to 0? {[}1, 2, 3{]} has numbers >= 0, but result is False. So that's not it. Let's re{-}examine. ({[}2, 1{]}, 2) {-}> True. ({[}1, 2, 3{]}, 9) {-}> False. ({[}0{]}, 0) {-}> True. ({[}5, {-}5, 5{]}, {-}51) {-}> False. ({[}{-}10{]},} \\
            }
        }
    }
    & & \\ \\

    \theutterance \stepcounter{utterance}  
    & & & \multicolumn{4}{p{0.6\linewidth}}{
        \cellcolor[rgb]{0.9,0.9,0.9}{
            \makecell[{{p{\linewidth}}}]{
                \texttt{\tiny{[P1$\langle$GM]}}
                \texttt{Game Violation: ParseError: Response is not valid JSON.} \\
\texttt{Output ONLY:} \\
\texttt{One JSON with confidence\_score, current\_mode, input\_rationale, and hypothesis keys ...} \\
\texttt{OR} \\
\texttt{SOLVE: ```python ... ```} \\
            }
        }
    }
     \\ \\

    \theutterance \stepcounter{utterance}  
    & \multicolumn{4}{p{0.6\linewidth}}{
        \cellcolor[rgb]{0.561,0.737,0.561}{
            \makecell[{{p{\linewidth}}}]{
                \texttt{\tiny{[P1$\rangle$GM]}}
                \texttt{\{} \\
\texttt{  "current\_mode": "probing",} \\
\texttt{  "input\_rationale": "I have five examples: ({[}2, 1{]}, 2) {-}> True, ({[}1, 2, 3{]}, 9) {-}> False, ({[}0{]}, 0) {-}> True, ({[}5, {-}5, 5{]}, {-}51) {-}> False, ({[}{-}10{]}, 0) {-}> True. The 'contains' hypothesis is now inconsistent because {[}{-}10{]} does not contain 0, but the result is True. The 'sum equals' hypothesis is also inconsistent because the sum of {[}{-}10{]} is {-}10, not 0. Let's look at other properties. Maybe it checks if the list contains a non{-}negative number? {[}2, 1{]} has non{-}negatives, {[}1, 2, 3{]} has non{-}negatives but result is False. So that's not it. Maybe it checks if the list contains a number greater than or equal to 0? {[}1, 2, 3{]} has numbers >= 0, but result is False. So that's not it. Let's re{-}examine. ({[}2, 1{]}, 2) {-}> True. ({[}1, 2, 3{]}, 9) {-}> False. ({[}0{]}, 0) {-}> True. ({[}5, {-}5, 5{]}, {-}51) {-}> False. ({[}{-}10{]},} \\
            }
        }
    }
    & & \\ \\

    \theutterance \stepcounter{utterance}  
    & & & \multicolumn{4}{p{0.6\linewidth}}{
        \cellcolor[rgb]{0.9,0.9,0.9}{
            \makecell[{{p{\linewidth}}}]{
                \texttt{\tiny{[P1$\langle$GM]}}
                \texttt{Game Violation: ParseError: Response is not valid JSON.} \\
\texttt{Output ONLY:} \\
\texttt{One JSON with confidence\_score, current\_mode, input\_rationale, and hypothesis keys ...} \\
\texttt{OR} \\
\texttt{SOLVE: ```python ... ```} \\
            }
        }
    }
     \\ \\

    \theutterance \stepcounter{utterance}  
    & \multicolumn{4}{p{0.6\linewidth}}{
        \cellcolor[rgb]{0.561,0.737,0.561}{
            \makecell[{{p{\linewidth}}}]{
                \texttt{\tiny{[P1$\rangle$GM]}}
                \texttt{\{} \\
\texttt{  "current\_mode": "probing",} \\
\texttt{  "input\_rationale": "I have five examples: ({[}2, 1{]}, 2) {-}> True, ({[}1, 2, 3{]}, 9) {-}> False, ({[}0{]}, 0) {-}> True, ({[}5, {-}5, 5{]}, {-}51) {-}> False, ({[}{-}10{]}, 0) {-}> True. The 'contains' hypothesis is now inconsistent because {[}{-}10{]} does not contain 0, but the result is True. The 'sum equals' hypothesis is also inconsistent because the sum of {[}{-}10{]} is {-}10, not 0. Let's look at other properties. Maybe it checks if the list contains a non{-}negative number? {[}2, 1{]} has non{-}negatives, {[}1, 2, 3{]} has non{-}negatives but result is False. So that's not it. Maybe it checks if the list contains a number greater than or equal to 0? {[}1, 2, 3{]} has numbers >= 0, but result is False. So that's not it. Let's re{-}examine. ({[}2, 1{]}, 2) {-}> True. ({[}1, 2, 3{]}, 9) {-}> False. ({[}0{]}, 0) {-}> True. ({[}5, {-}5, 5{]}, {-}51) {-}> False. ({[}{-}10{]},} \\
            }
        }
    }
    & & \\ \\

    \theutterance \stepcounter{utterance}  
    & & & \multicolumn{2}{p{0.3\linewidth}}{
        \cellcolor[rgb]{0.9,0.9,0.9}{
            \makecell[{{p{\linewidth}}}]{
                \texttt{\tiny{[GM$|$GM]}}
                \texttt{Game Verdict: LOSS (turn limit reached)} \\
            }
        }
    }
    & & \\ \\

\end{supertabular}
}

\end{document}
